\section{Primera iteración de la Fase 6: El fondo y los sprites en el ESP32}
\label{sec:F6I01}

Esta iteración es de DISEÑO: no toca código. Deja documentado cómo se guarda la imagen del fondo (img\_f) en el ESP32 y cómo se presentan y borran los sprites encima de ella, de la forma más compacta posible y sin que el aparato tenga que procesar mucho. Se escribe antes de empezar el firmware para tener previsto todo (formato en el GWC, algoritmo del compositor, memoria, tiempos, pruebas y cambios en el Designer) y que el desarrollo posterior sea una aplicación directa de lo aquí recogido.

\subsection{Datos de partida}
\label{sec:F6I01Datos}

\begin{itemize}
	\item \textbf{Pantalla:} 320$\times$240 píxeles en el equipo de pruebas, pero el diseño debe valer para pantallas algo mayores (p.~ej. 480$\times$320) y, como máximo, para el rango que admite el GWC (320$\times$240 a 1024$\times$600). El controlador habitual (ILI9341, ST7789 u otro similar) se maneja por SPI y la memoria de vídeo está \textbf{en la pantalla, no en el ESP32}.
	\item \textbf{ESP32 sin PSRAM} en el equipo de pruebas. La RAM interna útil ronda los 320\,KB y, descontando pila, sistema y lo que ocupe el programa, la libre es bastante menor y fragmentada. Aunque el hardware final previsto (Waveshare ESP32-S3-LCD-2.8) pueda traer PSRAM, \textbf{el diseño no debe depender de ella}.
	\item \textbf{Escritura sola:} en la práctica la pantalla SPI se trata como de solo escritura; no se puede leer de ella lo que hay dibujado (o no compensa por lentitud).
	\item \textbf{Sprites monocromos en negro,} como en las Game \& Watch originales: no hay otros colores posibles. Cada sprite es una cadena de bits en la que 0 es <<fondo (transparente)>> y 1 es <<negro>>.
	\item \textbf{Posición fija:} un sprite no se mueve por píxeles; solo puede estar PRESENTADO en su sitio o BORRADO. Hay una tabla con el ancho, el alto y la posición (x, y) de cada uno, de modo que su rectángulo (caja) es constante y se conoce desde que se carga el cartucho.
	\item \textbf{Los sprites no se solapan entre sí en píxeles} (lo valida el Designer con la comparación por diferencia de píxeles). Atención: que no se solapen sus píxeles no implica que no se solapen sus cajas.
	\item \textbf{El fondo es estático:} se pinta al empezar y no cambia durante la partida. En el cartucho es opcional; si no hay img\_f vale el \texttt{ColorFondo} del Bloque 0 (F5-I13).
	\item El Motor Gráfico ya redibuja solo los segmentos cuya marca de cambio está activa (F1), de modo que el firmware recibe peticiones del tipo <<presentar el sprite $i$>> o <<borrar el sprite $i$>>.
\end{itemize}

\subsection{Análisis}
\label{sec:F6I01Analisis}

\subsubsection{Presupuesto de memoria}

La imagen de fondo ocupa lo que muestra la tabla~\ref{tab:F6I01Tamanos} según el formato. Los números son exactos (píxeles $\times$ bytes por píxel, más la paleta cuando existe).

\begin{table}[H]
	\centering
	\begin{tabular}{lrrr}
		\toprule
		Formato & 320$\times$240 & 480$\times$320 & 1024$\times$600 \\
		\midrule
		RGB565 crudo (2 bytes/píxel) & 153\,600 & 307\,200 & 1\,228\,800 \\
		Paleta de 8 bits (+512 de paleta) & 77\,312 & 154\,112 & 614\,912 \\
		Paleta de 4 bits (+32 de paleta) & 38\,432 & 76\,832 & 307\,232 \\
		\bottomrule
	\end{tabular}
	\caption{Tamaño del fondo en bytes según el formato y la resolución}
	\label{tab:F6I01Tamanos}
\end{table}

\begin{itemize}
	\item Un framebuffer completo RGB565 de 320$\times$240 son 150\,KB. Sin PSRAM es inviable en la práctica (no hay un bloque contiguo tan grande con el resto del programa funcionando) y a 480$\times$320 (300\,KB) es imposible. Por eso se descarta como solución general.
	\item La flash, en cambio, es abundante (4\,MB o más): 150\,KB o incluso 300\,KB son pocos. El cartucho no necesita comprimir el fondo para que quepa; comprimirlo solo ahorra flash, que no es el recurso escaso.
	\item \textbf{El recurso escaso es la RAM,} y es la que mide la iteración F5-I20 (MA02). Por tanto el fondo debe quedarse en la flash y leerse desde ahí sin copiarlo entero a la RAM.
\end{itemize}

\subsubsection{Ideas estudiadas y descartadas}

\begin{table}[H]
	\centering
	\begin{tabular}{p{4.2cm}p{9cm}}
		\toprule
		Idea & Valoración \\
		\midrule
		Dibujar con XOR (como en el simulador) & No sirve: el XOR solo restaura el fondo si este es de un único color conocido (blanco). Con una imagen de fondo cualquiera produce colores falsos. \\
		Guardar en RAM el trozo de fondo bajo cada sprite antes de presentarlo y devolverlo al borrar & Funciona, pero gasta RAM por sprite (y ya sabemos que la RAM es lo escaso) para guardar algo que \emph{ya está en la flash}. Con fondo estático, releer de la flash es mejor. \\
		Framebuffer completo en la RAM & Lo más sencillo, pero no cabe (ver arriba). Solo tendría sentido con PSRAM. \\
		Fondo comprimido en JPEG o PNG & Impide el acceso a un rectángulo cualquiera, exige descomprimir todo en la RAM, gasta CPU y el JPEG ensucia los colores planos. \\
		Leer la pantalla para restaurar & No es posible o es muy lento en los controladores SPI. \\
		\bottomrule
	\end{tabular}
	\caption{Alternativas descartadas}
	\label{tab:F6I01Descartadas}
\end{table}

\subsubsection{La idea elegida: la flash ya es la copia guardada}

La idea de <<guardar el trozo de fondo y volver a ponerlo al borrar>> se simplifica: el trozo de fondo que hay debajo de cualquier sprite \textbf{ya está guardado} en la flash, sin gastar nada de RAM. Borrar un sprite es volver a pintar su rectángulo con el fondo leído de la flash. Falta resolver una única cuestión, que es la del siguiente punto.

\subsubsection{El problema de las cajas que se solapan}

Un sprite en diagonal (una pierna, un brazo, una cuerda) tiene una caja grande en la que sus píxeles negros son pocos. Esa caja puede contener a su vez píxeles de OTRO sprite que esté presentado. Si para borrar A se repinta toda su caja con el fondo, se borra un trozo de B, que sigue presentado. Ocurre igualmente al presentar: la caja de A pisaría con fondo parte de B si se pintase la caja entera.

La solución es el \textbf{compositor de rectángulos sucios:} al cambiar un sprite no se pinta su máscara sino que se \emph{compone} su caja entera en un búfer pequeño de RAM (fondo leído de la flash más los sprites presentados que cortan esa caja) y se envía ese búfer a la pantalla de una vez. Como los píxeles de los sprites no se solapan, el orden de composición no afecta al resultado.

Esto trae además dos ventajas: \textbf{no hay parpadeo} (nunca se ve el borrado y después el dibujado, la pantalla recibe el resultado final) y no hace falta ninguna RAM proporcional al número de sprites.

\subsection{Diseño}
\label{sec:F6I01Diseno}

\subsubsection{Modelo de datos en el firmware}

\begin{itemize}
	\item \textbf{Tabla de sprites} (constante, en la flash, sale del GWC; los Bloques 2 y 3 ya contienen estos datos): por cada sprite \texttt{ancho}, \texttt{alto}, \texttt{x}, \texttt{y} y el desplazamiento de su máscara. La caja es $[x,\,x+ancho)\times[y,\,y+alto)$.
	\item \textbf{Máscara:} cadena de bits continua, de $ancho\times alto$ bits, por filas, de izquierda a derecha y de arriba a abajo; el bit más significativo del byte es el primero. No se rellena cada fila hasta el byte: así el formato coincide con el que ya genera el GWC, y el compositor lleva un índice de bit que simplemente se incrementa (sin multiplicaciones por píxel).
	\item \textbf{Estado de presentación} (RAM, 1 bit por sprite): \texttt{visible[i]}. Es todo lo que el compositor necesita saber en tiempo de ejecución.
	\item \textbf{Lista de vecinos} de cada sprite: los otros sprites cuya caja corta con la suya. Se calcula una sola vez al cargar el cartucho (con 54 sprites son $54^2\approx 2\,900$ comparaciones de rectángulos, despreciable) y se guarda como una lista corta de índices. Se puede calcular en el firmware o en el Designer; se decide al implementar (ver \ref{sec:F6I01Abiertos}).
	\item Un sprite sin vecinos (lo habitual) usa el camino rápido descrito más abajo.
\end{itemize}

\subsubsection{Formato del fondo en el GWC (Bloque 5)}

Hoy el Bloque 5 del GWC está reservado y vacío (F5-I06). Se propone el siguiente contenido, que mantiene el criterio de que el ESP32 no procese ni interprete nada al leerlo:

\begin{table}[H]
	\centering
	\begin{tabular}{llp{7.5cm}}
		\toprule
		Campo & Tamaño & Contenido \\
		\midrule
		\texttt{formato} & 1 byte & 0 = RGB565 crudo; 1 = paleta de 8 bits; 2 = paleta de 4 bits (reservados 3 en adelante: teselas, compresión) \\
		\texttt{ancho}, \texttt{alto} & 2 + 2 bytes & Deben coincidir con la resolución del Setup (el Designer ya escala img\_f a esa resolución, F5-I13) \\
		\texttt{numColores} & 2 bytes & Solo para paleta; 0 en RGB565 \\
		paleta & $numColores\times 2$ & Colores ya en RGB565 y en el orden de bytes de la pantalla \\
		datos & según formato & Píxeles por filas, sin relleno entre filas salvo en 4 bits (el último nibble de una fila impar se rellena para que cada fila empiece en byte) \\
		\bottomrule
	\end{tabular}
	\caption{Contenido propuesto para el Bloque 5 (fondo)}
	\label{tab:F6I01Bloque5}
\end{table}

\begin{itemize}
	\item \textbf{RGB565 con los bytes ya en el orden de la pantalla} (normalmente byte alto primero). Así el firmware copia filas de la flash a un búfer y las envía tal cual, sin intercambiar bytes píxel a píxel. El orden depende del controlador, por lo que se fija en una cabecera del cartucho o en una constante del firmware, y se prueba en el equipo real.
	\item \textbf{Formato por defecto: RGB565 crudo.} Es lo más simple y no cuesta nada de CPU. Con 150\,KB a 320$\times$240 está sobrado.
	\item \textbf{Paleta de 4 u 8 bits:} el Designer cuenta los colores exactos de img\_f; si hay 16 o menos usa 4 bits, si hay 256 o menos usa 8 bits, y si hay más usa RGB565 crudo (no se cuantiza ni se hace dithering sin que el usuario lo pida). Cuesta una consulta a una tabla de 16 o 256 entradas por píxel, que es prácticamente gratis. Atención a la lección de F5 (el antialiasing de los editores de imagen convierte unos pocos colores en cientos): un fondo con muchos colores <<casi iguales>> no se reducirá a paleta y se quedará en RGB565.
	\item \textbf{Sin img\_f:} el Bloque 5 va vacío y se usa \texttt{ColorFondo} (Bloque 0). El compositor entonces no lee la flash: rellena el búfer con ese color. Ver también el arranque.
	\item \textbf{Fondo en la flash:} el GWC (o su Bloque 5) se guarda en una partición de datos. Se accede por memoria mapeada (\texttt{esp\_partition\_mmap}), con lo que leer una fila del fondo es leer un bloque contiguo de memoria; la caché de la flash hace el resto. Si se prefiere no mapear, \texttt{esp\_partition\_read} copia una fila al búfer con el mismo resultado un poco más lento. El mapeo se hace de la partición entera y se accede por desplazamiento, ya que el mapeo exige alineación de 64\,KB.
\end{itemize}

\subsubsection{Algoritmo del compositor}

El Motor Gráfico no pinta directamente: actualiza \texttt{visible[i]} y marca los sprites cambiados. Al terminar la iteración se vacía la lista de cambios con el algoritmo siguiente. \textbf{Primero se actualizan todos los estados y después se compone,} de otro modo la composición de un sprite vería un estado intermedio de otro que cambia en la misma iteración.

\begin{verbatim}
	para cada sprite i marcado como cambiado:
	caja = (x[i], y[i], ancho[i], alto[i])
	si vecinos[i] está vacío:                 // camino rápido
	si visible[i]:  copiar filas del fondo + poner negro la máscara
	si no:          copiar filas del fondo (nada más)
	si no:                                    // camino general
	copiar filas del fondo
	para cada sprite j en {i} + vecinos[i]:
	si visible[j]: poner negro la máscara de j, recortada a la caja
	enviar el búfer a la pantalla (ventana = caja)
	quitar la marca de i
\end{verbatim}

\begin{itemize}
	\item \textbf{Copiar filas del fondo:} para cada fila de la caja se copia \texttt{ancho[i]} píxeles desde \texttt{fondo + (fila\,$\times$\,anchoPantalla + x)\,$\times$\,2} (en RGB565) al búfer. Con paleta se convierte con la tabla de colores al copiar.
	\item \textbf{Poner negro la máscara:} se recorre la cadena de bits de la máscara; por cada 1 se escribe \texttt{0x0000} (negro en RGB565, que además es el mismo en cualquier orden de bytes) en la posición correspondiente del búfer. Un byte de la máscara a 0 se salta de golpe (los sprites suelen tener mucho vacío). Con un vecino, se recorta a la parte de su máscara que cae dentro de la caja.
	\item \textbf{Dos sprites cambiados con cajas que se tocan en la misma iteración:} cada uno compone su caja con el estado ya final de todos, de modo que el resultado es correcto y solo se repinta un par de veces la zona común. No hace falta fusionar rectángulos.
	\item \textbf{Cajas grandes:} el búfer tiene un tamaño fijo y se procesa la caja por bandas de $N$ filas (p.~ej. 16). Con $N=16$ y una caja de la anchura de la pantalla (320) el búfer es de $320\times16\times2=10\,240$ bytes; con dos búferes (uno se envía mientras se compone el otro) son 20\,KB. Es un presupuesto fijo e independiente del número de sprites.
	\item \textbf{Coste del camino rápido al borrar:} es copiar filas de la flash a un búfer y enviarlas; sin ninguna lógica de máscaras.
\end{itemize}

\subsubsection{Envío a la pantalla}

\begin{itemize}
	\item Cada rectángulo se envía como: fijar ventana (columnas y filas), orden de escritura de memoria y los datos. El firmware se limita a dos funciones propias, \texttt{pantallaVentana(x, y, w, h)} y \texttt{pantallaEnviar(búfer, bytes)}, de modo que cambiar de controlador o de biblioteca (esp\_lcd, TFT\_eSPI, LovyanGFX, \ldots) no toque el compositor.
	\item \textbf{DMA:} el búfer debe estar en la RAM interna (no en la flash) y alineado. El DMA del SPI no puede leer directamente de la flash mapeada, y por eso se copia primero al búfer. Con dos búferes se solapa la composición con el envío.
	\item \textbf{Reloj del SPI:} 40\,MHz como punto de partida; muchos paneles admiten más (hasta 80\,MHz) y se sube después de medir.
\end{itemize}

\subsubsection{Arranque y cambio de partida}

\begin{itemize}
	\item Al arrancar (o al volver a empezar) se pinta el fondo completo por bandas, copiando filas de la flash al búfer y enviándolas. Después se pintan los sprites que arrancan presentados, usando el compositor normal.
	\item Si no hay img\_f, el fondo es un color liso: se rellena el búfer una sola vez con \texttt{ColorFondo} y se reutiliza en todas las bandas, sin tocar la flash.
	\item Si el cartucho cambia o se carga otro, se vuelve a ejecutar el arranque. No hay estado del fondo que invalidar porque el fondo nunca se guarda en la RAM.
\end{itemize}

\subsection{Estimaciones de tiempo}
\label{sec:F6I01Tiempos}

Son órdenes de magnitud para decidir; se medirán en el equipo (\ref{sec:F6I01Pruebas}).

\begin{table}[H]
	\centering
	\begin{tabular}{lrr}
		\toprule
		Operación & Datos & Tiempo (SPI a 40\,MHz $\approx$ 5\,MB/s) \\
		\midrule
		Sprite de 20$\times$30 & 1\,200 bytes & $\approx 0'24$\,ms \\
		Caja de 100$\times$100 & 20\,000 bytes & $\approx 4$\,ms \\
		Fondo completo 320$\times$240 & 153\,600 bytes & $\approx 31$\,ms \\
		Fondo completo 480$\times$320 & 307\,200 bytes & $\approx 61$\,ms \\
		\bottomrule
	\end{tabular}
	\caption{Tiempo de envío por SPI según el tamaño del rectángulo}
	\label{tab:F6I01Tiempos}
\end{table}

\begin{itemize}
	\item Una iteración típica (decenas de sprites cambiando como mucho, casi siempre pocos) queda en pocos milisegundos. Las Game \& Watch iteran a decenas o centenas de milisegundos, con lo que hay mucho margen.
	\item El fondo completo solo se pinta al arrancar.
	\item La lectura de la flash mapeada y la composición añaden poco frente al SPI, pero se medirán por separado (\ref{sec:F6I01Pruebas}).
\end{itemize}

\subsection{Impacto en el resto del sistema}
\label{sec:F6I01Impacto}

\begin{itemize}
	\item \textbf{GWC (Bloque 5):} pasa de reservado a llevar el fondo con el formato anterior. El principio de F5-I13 (<<un GWC debe contener los mismos datos que el GWI y el juego debe funcionar igual>>) se mantiene: el lector del simulador debe poder interpretar el Bloque 5 y mostrar el mismo fondo. Como el RGB565 pierde bits de color, el simulador debería mostrar el fondo ya redondeado a RGB565 al cargar un GWC, para que lo que se ve en el navegador sea lo que verá el ESP32.
	\item \textbf{Designer:} un generador del Bloque 5 (conversión a RGB565 en el orden de bytes elegido, detección de 16 o 256 colores exactos, paleta, CRC), y la prueba de equivalencia GWI/GWC ampliada con el fondo.
	\item \textbf{F5-I20 (memoria):} la FLASH informada es el tamaño del bloque de fondo, que hasta ahora era siempre 0. En cuanto el Bloque 5 deje de ser nulo, el aviso del porcentaje lo recoge sin cambiar nada, como se dejó previsto. La RAM del GWC no cambia por el fondo.
	\item \textbf{RAM del firmware:} los búferes del compositor (unos 20\,KB), la tabla de sprites y de vecinos y el bit de visibilidad de cada sprite no los cuenta el GWC; se suman a la RAM que el usuario declara en el Setup para saber cuánto se deja libre.
\end{itemize}

\subsection{Pruebas previstas}
\label{sec:F6I01Pruebas}

\begin{itemize}
	\item \textbf{Compositor en el PC:} el compositor se escribe como una función pura (entrada: fondo, tabla de sprites, estados y caja; salida: búfer), sin dependencias del ESP32, de modo que se compile y se pruebe en el PC.
	\item \textbf{Equivalencia por propiedad:} tras una secuencia aleatoria de <<presentar/borrar>> componiendo solo las cajas sucias, la pantalla (simulada en un búfer del PC) debe ser idéntica a la que sale de pintar desde cero el fondo y todos los sprites presentados. Se repite con muchas secuencias y con cajas solapadas a propósito (el caso difícil).
	\item \textbf{Imagen de referencia:} se comparan las salidas con la imagen que genera el simulador JavaScript para el mismo cartucho (la del Parachute como referencia), una vez convertido el fondo a RGB565.
	\item \textbf{Formatos del fondo:} el mismo fondo en RGB565 crudo, paleta de 8 bits y paleta de 4 bits da la misma imagen compuesta.
	\item \textbf{Casos límite:} sprite en el borde de la pantalla, sprite de una sola fila o columna, dos sprites cuya caja se solapa y cuyas máscaras no, máscara vacía, fondo liso por \texttt{ColorFondo}, y cartucho con sprites fuera de la pantalla (se rechaza al cargar con un mensaje claro, igual que el resto de validaciones del lector).
	\item \textbf{Medidas en el ESP32:} microsegundos por iteración (lectura de la flash, composición y envío por separado), RAM libre mínima durante una partida y comprobación visual de que no hay parpadeo.
\end{itemize}

\subsection{Puntos abiertos}
\label{sec:F6I01Abiertos}

\begin{itemize}
	\item \textbf{Lista de vecinos:} calcularla en el firmware al cargar (más sencillo, un poco más de arranque) o guardarla en el GWC (más datos en el cartucho, nada que calcular). Mi inclinación es calcularla en el firmware.
	\item \textbf{Orden de bytes y controlador:} depende de la pantalla definitiva (ILI9341, ST7789, \ldots); se comprueba en el equipo real y, de hacer falta, se parametriza.
	\item \textbf{Pantallas mayores:} más allá de 480$\times$320 el SPI se queda corto para repintar el fondo entero (60\,ms o más), aunque los sprites siguen siendo baratos. Es aceptable porque el fondo solo se pinta al empezar.
	\item \textbf{Pantalla que sí permita leer, o hardware con PSRAM:} no se usan. El diseño funciona sin ellos y sin cambios; si más adelante se quisiera un framebuffer completo en el S3 con PSRAM sería una variante del compositor, no un cambio de formato.
	\item \textbf{Teselas con deduplicación y RLE:} una optimización posterior, solo si el tamaño del cartucho llegara a importar. El campo \texttt{formato} del Bloque 5 deja sitio para ello sin cambiar nada más.
	\item \textbf{Ajuste del reloj del SPI y doble búfer:} se afinan con las medidas del equipo.
\end{itemize}

\subsection{Criterios de aceptación}
\label{sec:F6I01Criterios}

\begin{itemize}
	\item El fondo no se guarda en la RAM en ningún formato; solo se leen de la flash las filas que hacen falta.
	\item Presentar o borrar cualquier sprite deja la pantalla idéntica a la de pintar todo desde cero, incluso con cajas solapadas.
	\item La RAM usada por el compositor es fija y no depende del número de sprites.
	\item El fondo en los tres formatos da la misma imagen y el simulador muestra el mismo fondo que el ESP32.
	\item Sin parpadeo visible al presentar o borrar.
	\item Una iteración típica se resuelve en unos pocos milisegundos.
\end{itemize}